\begin{table}[!htbp]
\centering
\caption{Pre-Treatment Balance of Fiscal Fundamentals}
\label{tab:post_northeast_did_balance}
\small
\begin{tabular}{lcc}
\toprule
 & Success vs. & Failure vs. \\
 & no override & no override \\
\midrule
Population size (log) & -0.095 & -0.015 \\
 & (0.067) & (0.065) \\
Outstanding debt level (\%) & 4.269 & 1.375 \\
 & (3.982) & (5.021) \\
Unemployment rate (\%) & -0.356** & 0.031 \\
 & (0.153) & (0.191) \\
Revenue stability (\%) & 7.301*** & 2.450 \\
 & (1.240) & (1.683) \\
Government revenue (per capita) & 0.169 & -0.096 \\
 & (0.108) & (0.122) \\
Excess property tax capacity & -0.031*** & -0.048*** \\
 & (0.009) & (0.007) \\
Fiscal reserve (\%) & -2.073** & -2.573* \\
 & (0.983) & (1.402) \\
Budget balance (\%) & -0.722 & -1.105* \\
 & (0.640) & (0.653) \\
\bottomrule
\end{tabular}
\vspace{0.05in}
\begin{minipage}{0.96\textwidth}\footnotesize\emph{Notes:} Each cell is the coefficient from regressing the listed control (measured at the election fiscal year, $t-1$) on the treatment indicator within stacks, i.e.\ the pre-treatment difference between treated events and their no-override controls; standard errors in parentheses are clustered by municipality. A significant coefficient indicates pre-treatment imbalance (selection). Override municipalities of both types are lower on excess property tax capacity and fiscal reserve; successful-override municipalities are additionally higher on revenue stability and lower on unemployment. The two imbalanced reserve/capacity measures are also the controls that carry signal in the rating-change models (Tables~\ref{tab:post_northeast_did_downgrade} and~\ref{tab:post_northeast_did_upgrade}). * p $<$ 0.10, ** p $<$ 0.05, *** p $<$ 0.01.\end{minipage}
\end{table}
